\documentclass[../../../include/open-logic-section]{subfiles}

\begin{document}

\olfileid[hi]{sth}{story}{rus}
\olsection{रसेल का विरोधाभास (फिर से)}

\olref[sfr][][]{part} में हमने सहज समुच्चय सिद्धांत के साथ काम किया था।
किंतु एक \emph{अत्यंत} सहज अवधारणा के अनुसार समुच्चय केवल विधेयों के
विस्तार होते हैं। यह सहज विचार निम्न सिद्धांत को अनिवार्य बनाएगा:

\begin{defish}
  \emph{सहज समुच्चय-निर्माण।} प्रत्येक सूत्र $\phi$ के लिए
  $\Setabs{x}{\phi(x)}$ का अस्तित्व है।
\end{defish}

यह सिद्धांत जितना भी आकर्षक हो, इसकी असंगति सिद्ध की जा सकती है। हमने इसे
\olref[sfr][set][rus]{sec} में देखा था, किंतु परिणाम इतना महत्त्वपूर्ण और
इतना सीधा है कि उसे शब्दशः दोहराना उचित है।

\begin{thm}[रसेल का विरोधाभास]
ऐसा कोई समुच्चय $R = \Setabs{x}{x \notin x}$ नहीं है।
\end{thm}

\begin{proof}
यदि $R = \Setabs{x}{x \notin x}$ का अस्तित्व हो, तो
$R \in R$ तभी और केवल तभी जब $R \notin R$; यह विरोधाभास है।
\end{proof}

रसेल ने यह परिणाम जून 1901 में खोजा। (हालाँकि उन्होंने विरोधाभास को ठीक
उसी रूप में नहीं रखा था जिसमें हमने अभी प्रस्तुत किया है, क्योंकि वे
\emph{ग्रुंडगेज़ेत्से} में प्रतिपादित फ़्रेगे के समुच्चय सिद्धांत पर विचार कर
रहे थे। इस पर हम \olref[blv]{sec} में लौटेंगे।) रसेल ने 16 जून 1902 को
फ़्रेगे को पत्र लिखकर फ़्रेगे के तंत्र की असंगति समझाई। पत्राचार और कुछ
पृष्ठभूमि के लिए \citet[pp.~124--8]{Heijenoort1967} देखें।

इस बात पर बल देना उचित है कि यह दो-पंक्ति का प्रमाण \emph{शुद्ध तर्कशास्त्र}
का परिणाम है। माना कि हमने अपने संकेतन $\Setabs{x}{x \notin x}$ में एक
(अतार्किक?)\ अभिगृहीत, विस्तारिता, का निहित रूप से उपयोग किया; क्योंकि यदि
$\phi$-वस्तुओं का कोई समुच्चय हो, तो $\Setabs{x}{\phi(x)}$ को (विस्तारिता
के कारण) उनका \emph{एकमात्र} समुच्चय होना है। किंतु परिणाम को इस रूप में
कहकर हम विस्तारिता के संकेत तक से बच सकते हैं: \emph{ऐसा कोई समुच्चय नहीं
है जिसके सदस्य ठीक वे समुच्चय हों जो स्वयं के सदस्य नहीं हैं}। और इसका
समुच्चयों से कोई विशेष संबंध नहीं। स्वयं रसेल ने देखा था कि ठीक वैसा ही
तर्क आपको यह निष्कर्ष देगा: \emph{ऐसा कोई पुरुष नहीं जो ठीक उन पुरुषों की
हजामत बनाए जो अपनी हजामत स्वयं नहीं बनाते}। अथवा: \emph{ऐसा कोई पग कुत्ता
नहीं जो ठीक उन पग कुत्तों को सूँघे जो स्वयं को नहीं सूँघते}। और इसी तरह।
योजनाबद्ध रूप में परिणाम का आकार केवल यह है:
\[
\lnot \exists x \forall z(Rzx \liff \lnot R zz).
\]
और यह प्रथम-क्रम तर्कशास्त्र का केवल एक प्रमेय (योजना) है। फलतः अपने
समुच्चय सिद्धांत में थोड़ी फेरबदल करके हम रसेल के विरोधाभास से नहीं बच सकते;
यह तो समुच्चय सिद्धांत तक \emph{पहुँचने} से पहले ही उत्पन्न हो जाता है। यदि
हमें (चिरसम्मत) प्रथम-क्रम तर्कशास्त्र का उपयोग करना है, तो हमें बस यह
\emph{स्वीकार} करना होगा कि $R = \Setabs{x}{x\notin x}$ जैसा कोई समुच्चय
नहीं है।

निष्कर्ष यह है। यदि आप असंगति से \emph{बचते हुए} सहज समुच्चय-निर्माण को
स्वीकार करना चाहते हैं, तो केवल \emph{समुच्चय सिद्धांत} में फेरबदल पर्याप्त
नहीं होगा। इसके बजाय आपको अपने \emph{तर्कशास्त्र} का आमूल परिवर्तन करना
होगा।

निश्चय ही अचिरसम्मत तर्कशास्त्रों वाले समुच्चय सिद्धांत प्रस्तुत किए गए हैं।
पर वे---कम-से-कम इतना तो कहा ही जा सकता है---अमानक हैं। रसेल के विरोधाभास
के प्रति मानक दृष्टिकोण उसे अनस्तित्व का सीधा प्रमाण मानना और फिर उसके साथ
काम करना सीखना है। हम इसी दृष्टिकोण का अनुसरण करेंगे।

\end{document}
